% Part: incompleteness
% Chapter: representability-in-q
% Section: basic-representable

\documentclass[../../../include/open-logic-section]{subfiles}

\begin{document}

\olfileid{inc}{req}{bre}
\olsection{基础函数在~$\Th{Q}$ 中可表示}

首先我们需要证明所有基础函数在~$\Th{Q}$ 中都是可表示的。
归根结底，我们需要说明如何为每个 $k$ 元基础函数 $f(x_0,\dots,x_{k-1})$ 指派一个表示它的公式
$!A_f(x_0,\dots,x_{k-1},y)$。

我们将能表示零函数、后继函数、加法、乘法、等词的特征函数以及投影函数。在每种情况下，恰当的表示公式都十分直接；例如，零函数用公式 $y = \Obj 0$ 表示，后继函数用公式 $x_0' = y$ 表示，加法用公式 $(x_0 + x_1) = y$ 表示。具体工作在于证明 $\Th{Q}$ 能够证明相关的语句；例如，要说明加法由上述公式表示，就需要证明对于任意一对自然数 $m$ 和 $n$，$\Th{Q}$ 都能证明
\begin{align*}
& \eq[\num n + \num m][\num {n+m}] \text{ 且}\\
& \lforall[y][(\eq[(\num n + \num m)][y] \lif \eq[y][\num{n+m}])].
\end{align*}

\begin{prop}
\ollabel{prop:rep-zero}
零函数 $\Zero(x) = 0$ 在~$\Th{Q}$ 中由
$!A_{\Zero}(x,y) \ident \eq[y][\Obj 0]$ 表示。
\end{prop}

\begin{prop}
\ollabel{prop:rep-succ}
后继函数 $\Succ(x) = x+1$ 在~$\Th{Q}$ 中由
$!A_{\Succ}(x,y) \ident \eq[y][x']$ 表示。  
\end{prop}

\begin{prop}
\ollabel{prop:rep-proj}
投影函数 $\Proj{n}{i}(x_0, \dots, x_{n-1}) = x_i$ 在~$\Th{Q}$ 中由 \[!A_{\Proj{n}{i}}(x_0, \dots, x_{n-1}, y) \ident \eq[y][x_i].\] 表示。
\end{prop}

\begin{prob}
证明 $\eq[y][\Obj 0]$、$\eq[y][x']$ 和 $\eq[y][x_i]$ 分别表示
$\Zero$、$\Succ$ 和 $\Proj{n}{i}$。
\end{prob}

\begin{prop}
\ollabel{prop:rep-id}
等词的特征函数，
\[
\Char{=}(x_0, x_1) =
\begin{cases}
  1 & \text{ 若 } x_0 =x_1\\
  0 & otherwise
\end{cases}
\]，
在~$\Th{Q}$ 中由
\[
  !A_{\Char{=}}(x_0, x_1, y) \ident (\eq[x_0][x_1] \land \eq[y][\num{1}]) \lor (\eq/[x_0][x_1] \land
\eq[y][\num{0}]).
\] 表示。
\end{prop}

该证明需要以下引理。

\begin{lem}
\ollabel{lem:q-proves-neq} 给定自然数 $n$ 和 $m$，若 $n
\neq m$，则 $\Th{Q} \Proves \eq/[\num n][\num m]$。
\end{lem}

\begin{proof}
我们对 $n$ 进行归纳来证明：对于每个 $m$，若 $n \neq m$，则
$Q \Proves \eq/[\num n][\num m]$。

在基础情形中，$n = 0$。若 $m$ 不等于 $0$，则存在某个自然数 $k$ 使得 $m = k + 1$。
我们有一条公理表明
$\lforall[x][\eq/[0][x']]$。根据量词公理，将 $x$ 替换为
$\num k$，可以得出 $\eq/[0][\num k']$。而 $\num k'$ 就是
$\num m$。

在归纳步骤中，我们可以假设命题对 $n$ 成立，并考虑 $n+1$。
令 $m$ 为任意自然数。有两种可能：要么 $m = 0$，要么存在某个 $k$ 使得 $m = k+1$。
第一种情况的处理同上。在第二种情况下，假设 $n+1 \neq k+1$。
那么 $n \neq k$。根据关于 $n$ 的归纳假设，我们有
$\Th{Q} \Proves \eq/[\num n][\num k]$。我们有一条公理表明
$\lforall[x][\lforall[y][\eq[x'][y'] \lif \eq[x][y]]]$。根据量词公理，我们有
$\eq[\num n'][\num k'] \lif \eq[\num n][\num
  k]$。根据命题逻辑，我们在 $\Th{Q}$ 中可推出
$\eq/[\num n][\num k] \lif \eq/[\num n'][\num k']$。根据肯定前件律，我们可以推出
$\eq/[\num n'][\num k']$，这正是我们想要的结果，因为 $\num k'$ 就是 $\num m$。
\end{proof}

\begin{explain}
注意该引理并没有说出太多东西：本质上它说的是 $\Th{Q}$ 能证明不同的数字项指称不同的对象。例如，
$\Th{Q}$ 能证明 $0'' \neq 0'''$。然而，要证明这在一般情况下也成立，则需要格外小心。另外还要注意，虽然我们使用了归纳法，但这是 $\Th{Q}$ \emph{之外}的归纳法。
\end{explain}

\begin{proof}[证明 \olref{prop:rep-id}]
若 $n = m$，则 $\num{n}$ 和 $\num{m}$ 是同一个项，且
$\Char{=}(n, m) = 1$。而 $\Th{Q} \Proves (\eq[\num{n}][\num{m}] \land
\eq[\num{1}][\num{1}])$，因此可推出 $!A_=(\num{n}, \num{m},
\num{1})$。若 $n \neq m$，则 $\Char=(n, m) = 0$。根据
\olref{lem:q-proves-neq}，$\Th{Q} \Proves \eq/[\num{n}][\num{m}]$，因而也有
$(\eq/[\num{n}][\num{m}] \land \Obj 0 = \Obj 0)$。于是 $\Th{Q}
\Proves !A_=(\num{n}, \num{m}, \num{0})$。

对于第二部分，我们同样分两种情况。若 $n = m$，我们需要
证明 $\Th{Q} \Proves \lforall[y][(!A_=(\num{n}, \num{m}, y)
  \lif \eq[y][\num{1}])]$。非形式地论证如下：假设
$!A_=(\num{n},
\num{m}, y)$，即
\[
(\eq[\num{n}][\num{n}] \land \eq[y][\num{1}]) \lor
(\eq/[\num{n}][\num{n}] \land \eq[y][\num{0}])
\]
由左析取支根据逻辑可推出 $\eq[y][\num{1}]$；右析取支与
$\eq[\num{n}][\num{n}]$ 矛盾，而后者是逻辑可证的。

另一方面，假设 $n \neq m$。那么 $!A_=(\num{n},
\num{m}, y)$ 为
\[
(\eq[\num{n}][\num{m}] \land \eq[y][\num{1}]) \lor
(\eq/[\num{n}][\num{m}] \land \eq[y][\num{0}])
\]
这里，左析取支与 $\eq/[\num{n}][\num{m}]$ 矛盾，而根据
\olref{lem:q-proves-neq}，该式在 $\Th{Q}$ 中可证；右析取支蕴涵 $\eq[y][\num{0}]$。
\end{proof}

\begin{prop}
\ollabel{prop:rep-add}
加法函数 $\Add(x_0, x_1) = x_0+x_1$ 在~$\Th{Q}$ 中由
\[
  !A_{\Add}(x_0, x_1, y) \ident \eq[y][(x_0 + x_1)].
\] 表示。
\end{prop}

\begin{lem}
\ollabel{lem:q-proves-add}
$\Th{Q} \Proves \eq[(\num{n} + \num{m})][\num{n+m}]$
\end{lem}

\begin{proof}
我们通过对 $m$ 进行归纳来证明此结论。若 $m = 0$，结论是
$\Th{Q} \Proves \eq[(\num{n} + \Obj 0)][\num{n}]$。这由
公理 $!Q_4$ 即得。现在假设对 $m$ 的结论成立；我们来证明对 $m+1$ 的结论，
即证明 $\Th{Q} \Proves \eq[(\num{n} +
  \num{m+1})][\num{n+m+1}]$。注意 $\num{m+1}$ 就是 $\num{m}'$，
而 $\num{n+m+1}$ 就是 $\num{n+m}'$。根据公理 $!Q_5$，$\Th{Q}
\Proves \eq[(\num{n} + \num{m}')][(\num{n}+\num{m})']$。根据归纳
假设，$\Th{Q} \Proves \eq[(\num{n} + \num{m})][\num{n+m}]$。所以
$\Th{Q} \Proves \eq[(\num{n} + \num{m}')][\num{n+m}']$。
\end{proof}

\begin{proof}[\olref{prop:rep-add} 的证明]
表示 $\Add$ 的公式~$!A_\Add(x_0, x_1, y)$ 是
$\eq[y][(x_0 + x_1)]$。首先我们证明，如果 $\Add(n, m) = k$，那么
$\Th{Q} \Proves !A_\Add(\num{n}, \num{m}, \num{k})$，即 $\Th{Q}
\Proves \eq[\num{k}][(\num{n} + \num{m})]$。但由于 $k = n + m$，
$\num{k}$ 就是 $\num{n+m}$，而且我们已经在
\olref{lem:q-proves-add} 中证明了 $\Th{Q} \Proves \eq[(\num{n} +
  \num{m})][\num{n+m}]$。

我们还须证明，如果 $\Add(n, m) = k$，则
\[
\Th{Q} \Proves \lforall[y][(!A_\Add(\num{n}, \num{m}, y) \lif
  \eq[y][\num{k}])].
\]
假设我们有 $\eq[(\num{n} + \num{m})][y]$。由于
\[
\Th{Q} \Proves \eq[(\num{n}+\num{m})][\num{n+m}],
\]，
我们可以将左边替换为 $\num{n+m}$，从而对任意~$y$ 得到
$\eq[\num{n+m}][y]$。
\end{proof}

\begin{prop}
\ollabel{prop:rep-mult}
乘法函数 $\Mult(x_0, x_1) = x_0 \cdot x_1$ 在~$\Th{Q}$ 中由
\[
  !A_{\Mult}(x_0, x_1, y) \ident y = (x_0 \times x_1).
\] 表示。
\end{prop}

\begin{proof} 练习。\end{proof}

\begin{lem}
\ollabel{lem:q-proves-mult}
$\Th{Q} \Proves \eq[(\num{n} \times \num{m})][\num{n \cdot m}]$
\end{lem}

\begin{proof} 练习。\end{proof}

\begin{prob}
证明 \olref[inc][req][bre]{lem:q-proves-mult}。
\end{prob}

\begin{prob}
利用 \olref[inc][req][bre]{lem:q-proves-mult} 来证明
\olref[inc][req][bre]{prop:rep-mult}。
\end{prob}

\begin{explain}
  回想一下，我们用 $\times$ 表示算术语言的函数符号，而用 $\cdot$ 表示数上普通的乘法运算。因此 $\cdot$ 可以出现在表示数的表达式之间（例如在 $m \cdot n$ 中），而 $\times$ 仅出现在算术语言的项之间（例如在 $(\num{m} \times
  \num{n})$ 中）。更容易混淆的是，$+$ 既用作函数符号，也用作加法运算。当它出现在项之间时（例如在 $(\num{n} + \num{m})$ 中），它是算术语言的二元函数符号；当它出现在数之间时（例如在 $n+m$ 中），它是加法运算。这也包括 $\num{n+m}$ 这种情况：这是对应于数~$n+m$ 的标准数字项。
\end{explain}

\end{document}